% GENERATED FILE. Edit canonical lessons, then run npm run generate:books.
% canonical-source-hash: fnv1a64:282f8e646b0b8d38
% canonical-lessons: SA-C425-adbhutah, SA-C425-smarakam, SA-C425-preksakah, SA-C425-gavaksah, SA-C425-smaraniyah, SA-R425-review-418-425

\chapter{Monuments and Onlookers}
\label{ch:sa-monuments-and-onlookers}
\hlchaptermodality{425}

\begin{chapteropening}
\textbf{By the end of this chapter:} \emph{I can talk about a monument, the crowd and what is memorable.}

Complete the last lesson of chapter 425: I can talk about a monument, the crowd and what is memorable.
\end{chapteropening}

\section[\texorpdfstring{\sk{अद्भुतः} (adbhutaḥ)}{अद्भुतः (adbhutaḥ)}]{\sk{अद्भुतः} (adbhutaḥ) — marvellous; wonderful}

\label{lesson:SA-C425-adbhutah}



\begin{quote}
Before the new one: say the Sanskrit for a view, then the Sanskrit for strolls.
\end{quote}

\subsection*{You'll want to know: \sk{अद्भुतः}}
\textbf{\sk{अद्भुतः}} — \emph{adbhutaḥ} — \textquotedblleft{}marvellous; wonderful\textquotedblright{}.

\begin{quote}
\sk{मम} \sk{यात्रायाः} \sk{अनुभवः} \sk{अद्भुतः} \sk{अस्ति।}
\end{quote}

\begin{quote}
\emph{mama yātrāyāḥ anubhavaḥ adbhutaḥ asti.}
\end{quote}

\begin{quote}
The experience of my journey is marvellous.
\end{quote}

\begin{cousinweb}
\textbf{\sk{अद्भुतः}} (\emph{adbhutaḥ}) is the word MW gives for wonderful, marvellous; MW gives no sure derivation, so its origin is uncertain (MW).
\end{cousinweb}

\begin{grammarlens}[title={What you've built}]
One more word for getting around.
\end{grammarlens}

\subsection*{Your turn}
\begin{itemize}
\raggedright
  \item \emph{Say it:} \emph{adbhutaḥ}
  \item \emph{Say it:} \emph{adbhutaḥ}, once more
  \item \emph{Say it:} \emph{mama yātrāyāḥ anubhavaḥ adbhutaḥ asti.} — The experience of my journey is marvellous.
  \item \emph{Recall:} say the Sanskrit for a view, then the Sanskrit for strolls, then say \emph{adbhutaḥ} again
\end{itemize}

\begin{tcolorbox}[breakable,colback=teal!4,colframe=teal!35!black,title={Before you move on}]
What does \sk{अद्भुतः} mean? (\textquotedblleft{}Marvellous; wonderful\textquotedblright{}.) Say it once more.
\end{tcolorbox}
\section[\texorpdfstring{\sk{स्मारकम्} (smārakam)}{स्मारकम् (smārakam)}]{\sk{स्मारकम्} (smārakam) — a memorial}

\label{lesson:SA-C425-smarakam}



\begin{quote}
Before the new one: say the Sanskrit for strolls, then the Sanskrit for marvellous.
\end{quote}

\subsection*{You'll want to know: \sk{स्मारकम्}}
\textbf{\sk{स्मारकम्}} — \emph{smārakam} — \textquotedblleft{}a memorial\textquotedblright{}.

\begin{quote}
\sk{नगरस्य} \sk{मध्ये} \sk{एकं} \sk{स्मारकम्} \sk{अस्ति।}
\end{quote}

\begin{quote}
\emph{nagarasya madhye ekaṃ smārakam asti.}
\end{quote}

\begin{quote}
In the middle of the city there is a memorial.
\end{quote}

A neuter noun.

\begin{cousinweb}
\textbf{\sk{स्मारकम्}} (\emph{smārakam}) is built on the causative of the root √smṛ, \textquotedblleft{}to remember\textquotedblright{}, which also gives \emph{smarati}, \textquotedblleft{}remembers\textquotedblright{}: MW gives \emph{smāraka} as recalling, reminding of, and a memorial is its use as a noun in modern Sanskrit (MW).
\end{cousinweb}

\begin{grammarlens}[title={What you've built}]
One more word for getting around.
\end{grammarlens}

\subsection*{Your turn}
\begin{itemize}
\raggedright
  \item \emph{Say it:} \emph{smārakam}
  \item \emph{Say it:} \emph{smārakam}, once more
  \item \emph{Say it:} \emph{nagarasya madhye ekaṃ smārakam asti.} — In the middle of the city there is a memorial.
  \item \emph{Recall:} say the Sanskrit for strolls, then the Sanskrit for marvellous, then say \emph{smārakam} again
\end{itemize}

\begin{tcolorbox}[breakable,colback=teal!4,colframe=teal!35!black,title={Before you move on}]
What does \sk{स्मारकम्} mean? (\textquotedblleft{}A memorial\textquotedblright{}.) Say it once more.
\end{tcolorbox}
\section[\texorpdfstring{\sk{प्रेक्षकः} (prekṣakaḥ)}{प्रेक्षकः (prekṣakaḥ)}]{\sk{प्रेक्षकः} (prekṣakaḥ) — an onlooker; a spectator}

\label{lesson:SA-C425-preksakah}



\begin{quote}
Before the new one: say the Sanskrit for marvellous, then the Sanskrit for a memorial.
\end{quote}

\subsection*{You'll want to know: \sk{प्रेक्षकः}}
\textbf{\sk{प्रेक्षकः}} — \emph{prekṣakaḥ} — \textquotedblleft{}an onlooker; a spectator\textquotedblright{}.

\begin{quote}
\sk{स्मारकस्य} \sk{समीपे} \sk{प्रेक्षकाः} \sk{सन्ति।}
\end{quote}

\begin{quote}
\emph{smārakasya samīpe prekṣakāḥ santi.}
\end{quote}

\begin{quote}
There are onlookers near the memorial.
\end{quote}

A masculine noun.

\begin{cousinweb}
\textbf{\sk{प्रेक्षकः}} (\emph{prekṣakaḥ}) is built from \emph{pra}, \textquotedblleft{}forth\textquotedblright{}, and the root √īkṣ, \textquotedblleft{}to see\textquotedblright{}, with the agent ending \emph{-aka}; MW gives \emph{prekṣaka} as a spectator, a member of an audience (MW).
\end{cousinweb}

\begin{grammarlens}[title={What you've built}]
One more word for getting around.
\end{grammarlens}

\subsection*{Your turn}
\begin{itemize}
\raggedright
  \item \emph{Say it:} \emph{prekṣakaḥ}
  \item \emph{Say it:} \emph{prekṣakaḥ}, once more
  \item \emph{Say it:} \emph{smārakasya samīpe prekṣakāḥ santi.} — There are onlookers near the memorial.
  \item \emph{Recall:} say the Sanskrit for marvellous, then the Sanskrit for a memorial, then say \emph{prekṣakaḥ} again
\end{itemize}

\begin{tcolorbox}[breakable,colback=teal!4,colframe=teal!35!black,title={Before you move on}]
What does \sk{प्रेक्षकः} mean? (\textquotedblleft{}An onlooker; a spectator\textquotedblright{}.) Say it once more.
\end{tcolorbox}
\section[\texorpdfstring{\sk{गवाक्षः} (gavākṣaḥ)}{गवाक्षः (gavākṣaḥ)}]{\sk{गवाक्षः} (gavākṣaḥ) — a round window}

\label{lesson:SA-C425-gavaksah}



\begin{quote}
Before the new one: say the Sanskrit for a memorial, then the Sanskrit for an onlooker.
\end{quote}

\subsection*{You'll want to know: \sk{गवाक्षः}}
\textbf{\sk{गवाक्षः}} — \emph{gavākṣaḥ} — \textquotedblleft{}a round window\textquotedblright{}.

\begin{quote}
\sk{राज्ञी} \sk{गवाक्षात्} \sk{नगरं} \sk{पश्यति।}
\end{quote}

\begin{quote}
\emph{rājñī gavākṣāt nagaraṃ paśyati.}
\end{quote}

\begin{quote}
The queen looks out at the city from the round window.
\end{quote}

A masculine noun.

\begin{cousinweb}
\textbf{\sk{गवाक्षः}} (\emph{gavākṣaḥ}) is built from \emph{go}, \textquotedblleft{}a bull\textquotedblright{}, and \emph{akṣa}, \textquotedblleft{}an eye\textquotedblright{}: a bull's eye, as MW glosses it; MW gives \emph{gavākṣa} as an air-hole, a round window (MW).
\end{cousinweb}

\begin{grammarlens}[title={What you've built}]
One more word for getting around.
\end{grammarlens}

\subsection*{Your turn}
\begin{itemize}
\raggedright
  \item \emph{Say it:} \emph{gavākṣaḥ}
  \item \emph{Say it:} \emph{gavākṣaḥ}, once more
  \item \emph{Say it:} \emph{rājñī gavākṣāt nagaraṃ paśyati.} — The queen looks out at the city from the round window.
  \item \emph{Recall:} say the Sanskrit for a memorial, then the Sanskrit for an onlooker, then say \emph{gavākṣaḥ} again
\end{itemize}

\begin{tcolorbox}[breakable,colback=teal!4,colframe=teal!35!black,title={Before you move on}]
What does \sk{गवाक्षः} mean? (\textquotedblleft{}A round window\textquotedblright{}.) Say it once more.
\end{tcolorbox}
\section[\texorpdfstring{\sk{स्मरणीयः} (smaraṇīyaḥ)}{स्मरणीयः (smaraṇīyaḥ)}]{\sk{स्मरणीयः} (smaraṇīyaḥ) — memorable}

\label{lesson:SA-C425-smaraniyah}



\begin{quote}
Before the new one: say the Sanskrit for an onlooker, then the Sanskrit for a round window.
\end{quote}

\subsection*{You'll want to know: \sk{स्मरणीयः}}
\textbf{\sk{स्मरणीयः}} — \emph{smaraṇīyaḥ} — \textquotedblleft{}memorable\textquotedblright{}.

\begin{quote}
\sk{मम} \sk{तीर्थयात्रा} \sk{स्मरणीया} \sk{अस्ति।}
\end{quote}

\begin{quote}
\emph{mama tīrthayātrā smaraṇīyā asti.}
\end{quote}

\begin{quote}
My pilgrimage is memorable.
\end{quote}

\begin{cousinweb}
\textbf{\sk{स्मरणीयः}} (\emph{smaraṇīyaḥ}) is built on the root √smṛ, \textquotedblleft{}to remember\textquotedblright{}, which also gives \emph{smarati}, \textquotedblleft{}remembers\textquotedblright{}, with the ending \emph{-anīya}; MW gives \emph{smaraṇīya} as to be remembered, memorable (MW).
\end{cousinweb}

\begin{grammarlens}[title={What you've built}]
One more word for getting around.
\end{grammarlens}

\subsection*{Your turn}
\begin{itemize}
\raggedright
  \item \emph{Say it:} \emph{smaraṇīyaḥ}
  \item \emph{Say it:} \emph{smaraṇīyaḥ}, once more
  \item \emph{Say it:} \emph{mama tīrthayātrā smaraṇīyā asti.} — My pilgrimage is memorable.
  \item \emph{Recall:} say the Sanskrit for an onlooker, then the Sanskrit for a round window, then say \emph{smaraṇīyaḥ} again
\end{itemize}

\begin{tcolorbox}[breakable,colback=teal!4,colframe=teal!35!black,title={Before you move on}]
What does \sk{स्मरणीयः} mean? (\textquotedblleft{}Memorable\textquotedblright{}.) Say it once more.
\end{tcolorbox}
\section[(dialogue)]{Looking back: from Nets, Shells and Pearls to Monuments and Onlookers}

\label{lesson:SA-R425-review-418-425}



\begin{quote}
Say the Sanskrit for a net and a fishhook.
\end{quote}

\subsection*{Your turn}
Say these aloud:

\begin{itemize}
\raggedright
  \item a net, a fishhook, an oyster shell, a pearl, coral
  \item an oblong reservoir, a spring of water, a lotus pool, a water channel, a puddle
  \item a roadside water shed, gives a drink to, a mirage, comes upon, refreshes
  \item a hamlet, a suburb, a town square, a mansion, converses
  \item a town wall, a town gate, a moat, a tunnel, a ditch
  \item an entrance, an exit, opposite, a gap, a gatekeeper
  \item tours, worth seeing, delightful, a view, strolls
  \item marvellous, a memorial, an onlooker, a round window, memorable
\end{itemize}

\begin{tcolorbox}[breakable,colback=teal!4,colframe=teal!35!black,title={Before you move on}]
A net? (\textbf{\sk{जालम्}}, \emph{jālam}.) A fishhook? (\textbf{\sk{बडिशम्}}, \emph{baḍiśam}.) An oyster shell? (\textbf{\sk{शुक्तिः}}, \emph{śuktiḥ}.) A pearl? (\textbf{\sk{मौक्तिकम्}}, \emph{mauktikam}.) Coral? (\textbf{\sk{प्रवालम्}}, \emph{pravālam}.) An oblong reservoir? (\textbf{\sk{वापी}}, \emph{vāpī}.) A spring of water? (\textbf{\sk{उत्सः}}, \emph{utsaḥ}.) A lotus pool? (\textbf{\sk{पुष्करिणी}}, \emph{puṣkariṇī}.) A water channel? (\textbf{\sk{कुल्या}}, \emph{kulyā}.) A puddle? (\textbf{\sk{पल्वलम्}}, \emph{palvalam}.) A roadside water shed? (\textbf{\sk{प्रपा}}, \emph{prapā}.) Gives a drink to? (\textbf{\sk{पाययति}}, \emph{pāyayati}.) A mirage? (\textbf{\sk{मृगतृष्णा}}, \emph{mṛgatṛṣṇā}.) Comes upon? (\textbf{\sk{आसादयति}}, \emph{āsādayati}.) Refreshes? (\textbf{\sk{तर्पयति}}, \emph{tarpayati}.) A hamlet? (\textbf{\sk{पल्ली}}, \emph{pallī}.) A suburb? (\textbf{\sk{उपनगरम्}}, \emph{upanagaram}.) A town square? (\textbf{\sk{चत्वरम्}}, \emph{catvaram}.) A mansion? (\textbf{\sk{हर्म्यम्}}, \emph{harmyam}.) Converses? (\textbf{\sk{सम्भाषते}}, \emph{sambhāṣate}.) A town wall? (\textbf{\sk{प्राकारः}}, \emph{prākāraḥ}.) A town gate? (\textbf{\sk{गोपुरम्}}, \emph{gopuram}.) A moat? (\textbf{\sk{परिखा}}, \emph{parikhā}.) A tunnel? (\textbf{\sk{सुरंगा}}, \emph{suraṅgā}.) A ditch? (\textbf{\sk{खातम्}}, \emph{khātam}.) An entrance? (\textbf{\sk{प्रवेशः}}, \emph{praveśaḥ}.) An exit? (\textbf{\sk{निर्गमः}}, \emph{nirgamaḥ}.) Opposite? (\textbf{\sk{विपरीतः}}, \emph{viparītaḥ}.) A gap? (\textbf{\sk{अन्तरम्}}, \emph{antaram}.) A gatekeeper? (\textbf{\sk{द्वारपालः}}, \emph{dvārapālaḥ}.) Tours? (\textbf{\sk{पर्यटति}}, \emph{paryaṭati}.) Worth seeing? (\textbf{\sk{दर्शनीयः}}, \emph{darśanīyaḥ}.) Delightful? (\textbf{\sk{रमणीयः}}, \emph{ramaṇīyaḥ}.) A view? (\textbf{\sk{दृश्यम्}}, \emph{dṛśyam}.) Strolls? (\textbf{\sk{विहरति}}, \emph{viharati}.) Marvellous? (\textbf{\sk{अद्भुतः}}, \emph{adbhutaḥ}.) A memorial? (\textbf{\sk{स्मारकम्}}, \emph{smārakam}.) An onlooker? (\textbf{\sk{प्रेक्षकः}}, \emph{prekṣakaḥ}.) A round window? (\textbf{\sk{गवाक्षः}}, \emph{gavākṣaḥ}.) Memorable? (\textbf{\sk{स्मरणीयः}}, \emph{smaraṇīyaḥ}.)
\end{tcolorbox}
